The reaction subsystem combines kinetic laws, reaction-progress models,
reaction enthalpies, and reacting-material inventories into physical
thermal runaway reactions.

Each reaction contributes both a progress rate and a heat-generation
rate to the coupled system.


\subsection{Standard reaction model}

For a standard reaction, the progress rate is

\begin{equation}
    \dot{\alpha}
    =
    k(T)
    f
    \left(
        \alpha,
        \boldsymbol{\alpha}
    \right),
    \label{eq:reaction_progress}
\end{equation}

where $k(T)$ is the kinetic rate and $f$ is the reaction-progress
function.

If the reaction occupies an effective mass fraction $w$ of the total
cell mass and releases a specific enthalpy $\Delta H$, the
mass-specific heat-generation rate relative to total cell mass is

\begin{equation}
    \dot{q}
    =
    w
    \Delta H
    \dot{\alpha}.
    \label{eq:reaction_heat_specific}
\end{equation}

The total heat-generation rate for a cell of mass $m_{\mathrm{cell}}$
is therefore

\begin{equation}
    \dot{Q}
    =
    m_{\mathrm{cell}}
    w
    \Delta H
    \dot{\alpha}.
    \label{eq:reaction_heat_total}
\end{equation}


\subsection{Reaction enthalpy convention}

The parameter $\Delta H$ represents the effective heat release per unit
mass of reacting material.

For the reduced-order exothermic reactions currently implemented,
$\Delta H$ is treated as a positive heat-release quantity.

The resulting heat-generation term is therefore positive when

\begin{equation}
    \dot{\alpha} > 0.
\end{equation}

This convention avoids embedding thermodynamic sign conventions directly
into the thermal balance and is consistent with the current LiionTR
reaction implementation.


\subsection{Reaction mass fraction}

The parameter $w$ represents the effective fraction of total cell mass
participating in a particular reaction.

For literature models that report a volumetric reactant content
$W_r$ in \si{\kilogram\per\meter\cubed}, LiionTR converts this quantity
to a cell mass fraction using

\begin{equation}
    w
    =
    \frac{
        W_r
    }{
        \rho_{\mathrm{cell}}
    },
    \label{eq:volumetric_to_mass_fraction}
\end{equation}

where $\rho_{\mathrm{cell}}$ is the effective cell density.

This conversion enables literature models reported on a volumetric basis
to be used consistently within the cell-mass-based thermal formulation.


\subsection{Reaction context}

Individual reactions may depend on the state of other reactions.

LiionTR therefore constructs a reaction context containing the complete
reaction state

\begin{equation}
    \boldsymbol{\alpha}
    =
    \left[
        \alpha_1,
        \ldots,
        \alpha_N
    \right]^{\mathrm{T}}.
\end{equation}

The context provides access to quantities such as

\begin{equation}
    \alpha_j
\end{equation}

and

\begin{equation}
    1-\alpha_j,
\end{equation}

as well as transformed or normalized context variables.

This permits reaction coupling without storing mutable integration state
inside individual reaction objects.


\subsection{Multi-channel reactions}

Some literature models represent one physical decomposition process
using several parallel kinetic channels.

LiionTR represents these processes with a common conversion variable and
multiple kinetic channels.

The shared reaction progress rate is obtained from the sum of the
individual kinetic channels,

\begin{equation}
    \dot{\alpha}
    =
    f
    \left(
        \alpha,
        \boldsymbol{\alpha}
    \right)
    \sum_{j=1}^{N_c}
    k_j(T).
    \label{eq:multichannel_progress}
\end{equation}

For $N_c$ channels sharing a conversion $\alpha$, the heat-generation
rate can be written conceptually as


\begin{equation}
    \dot{q}
    =
    w
    f
    \left(
        \alpha,
        \boldsymbol{\alpha}
    \right)
    \sum_{j=1}^{N_c}
    \Delta H_j
    k_j(T).
    \label{eq:multichannel_heat}
\end{equation}

Each channel therefore contributes independently to heat generation,
while the physical reaction retains only one integrated progress state.


\subsection{Shared reaction progress}

The use of a shared conversion is physically distinct from representing
each kinetic channel as an independent reaction.

For a multi-channel decomposition mechanism, all channels consume the
same effective reactant inventory.

The model therefore integrates

\begin{equation}
    \alpha
\end{equation}

rather than introducing separate conversions

\begin{equation}
    \alpha_1,
    \alpha_2,
    \ldots,
    \alpha_{N_c}.
\end{equation}

This distinction is important for preserving a consistent reacting mass
inventory.


\subsection{Reaction network}

A thermal runaway mechanism is assembled as a network containing
multiple reactions.

For $N$ reactions, the conversion vector is

\begin{equation}
    \boldsymbol{\alpha}
    =
    \left[
        \alpha_1,
        \alpha_2,
        \ldots,
        \alpha_N
    \right]^{\mathrm{T}}.
\end{equation}

The corresponding progress-rate vector is

\begin{equation}
    \dot{\boldsymbol{\alpha}}
    =
    \left[
        \dot{\alpha}_1,
        \dot{\alpha}_2,
        \ldots,
        \dot{\alpha}_N
    \right]^{\mathrm{T}}.
\end{equation}

Each reaction is evaluated using the current temperature and current
reaction-network state.


\subsection{Network heat generation}

The total mass-specific heat-generation rate is

\begin{equation}
    \dot{q}_{\mathrm{total}}
    =
    \sum_{i=1}^{N}
    \dot{q}_i.
    \label{eq:network_specific_heat}
\end{equation}

The corresponding total cell heat-generation rate is

\begin{equation}
    \dot{Q}_{\mathrm{gen}}
    =
    m_{\mathrm{cell}}
    \dot{q}_{\mathrm{total}}.
    \label{eq:network_total_heat}
\end{equation}

This quantity is passed to the thermal model.


\subsection{Coupled thermal-reaction system}

The reaction network and thermal model form a strongly coupled
nonlinear system,

\begin{equation}
    \frac{dT}{dt}
    =
    \frac{
        m_{\mathrm{cell}}
        \displaystyle\sum_{i=1}^{N}
        \dot{q}_i
        -
        hA
        \left(
            T-T_{\infty}
        \right)
    }{
        C_{\mathrm{th}}
    },
    \label{eq:coupled_temperature}
\end{equation}

with

\begin{equation}
    \frac{d\alpha_i}{dt}
    =
    k_i(T)
    f_i
    \left(
        \alpha_i,
        \boldsymbol{\alpha}
    \right),
    \qquad
    i=1,\ldots,N.
    \label{eq:coupled_reactions}
\end{equation}

The reaction rates depend on temperature, while temperature depends on
reaction heat generation.

This bidirectional coupling produces the positive feedback characteristic
of thermal runaway.


\subsection{Initial reaction state}

Each reaction requires an initial conversion

\begin{equation}
    \alpha_{i,0}.
\end{equation}

The complete initial reaction state is

\begin{equation}
    \boldsymbol{\alpha}_0
    =
    \left[
        \alpha_{1,0},
        \ldots,
        \alpha_{N,0}
    \right]^{\mathrm{T}}.
\end{equation}

When a literature model reports an initial remaining fraction $r_0$,
LiionTR converts it using

\begin{equation}
    \alpha_0
    =
    1-r_0.
    \label{eq:remaining_to_conversion}
\end{equation}


\subsection{Software representation}

The physical organization of the reaction subsystem can be summarized
using the conceptual structure shown in
Figure~\ref{fig:reaction_architecture}.

\begin{figure}[htbp]
    \centering

    \begin{tikzpicture}[
        node distance=1.6cm and 2.0cm,
        block/.style={
            rectangle,
            rounded corners,
            draw,
            align=center,
            minimum width=3.1cm,
            minimum height=0.95cm
        },
        line/.style={
            -{Latex[length=3mm]},
            thick
        }
    ]

    \node[block] (kinetic) {Kinetic model\\$k(T)$};
    \node[block, right=of kinetic] (progress) {Progress model\\$f(\alpha,\mathbf{\alpha})$};

    \node[block, below=of $(kinetic)!0.5!(progress)$] (reaction)
        {Reaction};

    \node[block, below=of reaction] (network)
        {Reaction network};

    \node[block, below left=of network] (rates)
        {Progress rates\\$\dot{\boldsymbol{\alpha}}$};

    \node[block, below right=of network] (heat)
        {Heat generation\\$\dot{Q}_{\mathrm{gen}}$};

    \draw[line] (kinetic) -- (reaction);
    \draw[line] (progress) -- (reaction);
    \draw[line] (reaction) -- (network);
    \draw[line] (network) -- (rates);
    \draw[line] (network) -- (heat);

    \end{tikzpicture}

    \caption{
        Conceptual organization of the LiionTR reaction subsystem.
        Kinetic and progress models are combined into reaction models,
        which are assembled into a reaction network that supplies
        reaction progress rates and total heat generation.
    }
    \label{fig:reaction_architecture}
\end{figure}


\subsection{Current assumptions}

The current reaction framework assumes that:

\begin{itemize}
    \item every integrated reaction is represented by a scalar conversion;
    \item reaction states are spatially uniform;
    \item reaction enthalpies are effective constant quantities;
    \item reacting-material mass fractions remain fixed;
    \item pressure dependence is absent unless explicitly introduced by a
          custom model;
    \item reaction products are not automatically determined from
          chemical equilibrium;
    \item transport limitations are not explicitly resolved.
\end{itemize}


\subsection{Future extensions}

The modular reaction architecture is intended to support more detailed
physical descriptions.

Potential extensions include:

\begin{itemize}
    \item temperature-dependent reaction enthalpies;
    \item pressure-dependent kinetics;
    \item transport-limited reactions;
    \item competing reaction mechanisms;
    \item reaction-derived elemental source inventories;
    \item thermochemical equilibrium using Cantera;
    \item electrochemical-state-dependent kinetics;
    \item calibration against ARC and DSC measurements;
    \item parameter uncertainty and sensitivity analysis.
\end{itemize}